\documentclass[a4paper,11pt]{article}
\usepackage{exam-style-842-unit}
\examinehead{842 信息化概论 · 第 3 单元练习册}{练习二\quad 通信与网络技术}
\begin{document}

\unitcover{842 信息化概论 · 单元练习册}{第 3 单元\quad 计算机技术与网络技术}{练习二\quad 通信与网络技术（共 100 分）}{约 70 分钟}{通信技术的发展；信息传递方式与信息量；移动通信的四个阶段；计算机网络的定义与发展历程；三网融合的含义、层面与支撑技术；网络协议与网络体系结构；OSI 七层参考模型；TCP/IP 四层协议族；网络管理的概念、功能与 SNMP；网络安全的基本属性、技术机制与安全策略。}

\begin{knowmap}
\textbf{通信的本质：}快速、有效地实现信息的传递；从烽火狼烟、飞鸽传书 $\rightarrow$ 电报电话 $\rightarrow$ 移动通信 1G—4G $\rightarrow$ 卫星通信 $\rightarrow$ 全光通信与量子通信\\[2pt]
\textbf{计算机网络：}用通信线路把分散的独立计算机系统连接起来，按网络协议进行数据通信和资源共享；经历诞生 $\rightarrow$ 形成（ARPANET）$\rightarrow$ 互联互通（OSI、TCP/IP）$\rightarrow$ 高速网络（Internet）四个阶段\\[2pt]
\textbf{三网融合：}电信网·广播电视网·互联网 $\rightarrow$ 宽带通信网·数字电视网·下一代互联网；技术、业务、行业、终端、网络五个层面融合，不等于物理合一\\[2pt]
\textbf{协议与体系结构：}网络协议三要素＝语义·语法·同步；体系结构＝各层次及其协议的集合，分层带来独立性与标准化\\[2pt]
\textbf{两大体系：}OSI 七层（物理·数据链路·网络·运输·会话·表示·应用）与 TCP/IP 四层（网络接口·网络·传输·应用）\\[2pt]
\textbf{网络管理与安全：}管理站＋代理，SNMP 与 MIB；安全三属性＝机密性·完整性·可用性，技术机制＋管理措施并重
\end{knowmap}

\begin{keyterms}
\textbf{通信}：人们通过声音、文字、图像或动作相互沟通消息，本质是如何快速、有效地实现信息的传递。\\
\textbf{计算机网络}：用通信线路将分散在不同地点的独立的计算机系统相互连接，并按网络协议进行数据通信和实现资源共享的计算机集合。\\
\textbf{网络协议}：为进行计算机网络中的数据交换而建立的规则、标准或约定的集合，三要素为语义、语法和同步。\\
\textbf{网络体系结构}：计算机网络各层次及其协议的集合，分层使各层功能相对独立，便于实现、交流和标准化。\\
\textbf{OSI 参考模型}：国际标准化组织（ISO）制定的开放系统互连七层参考模型，是用于协调各层标准制定的概念性框架。\\
\textbf{TCP/IP}：Internet 使用的四层协议族，分为网络接口层、网络层、传输层和应用层，是事实上的工业标准。\\
\textbf{三网融合}：电信网、广播电视网、互联网在演进中技术功能趋于一致、业务范围趋于相同、互联互通、资源共享，主要指高层业务应用的融合。\\
\textbf{网络管理}：对网络资源进行监视、测试、配置、分析、评价和控制，使网络能有效、可靠、安全、经济地提供服务。\\
\textbf{SNMP}：简单网络管理协议，由互联网工程工作小组（IETF）定义，是基于 TCP/IP 协议族的网络管理标准。\\
\textbf{管理信息库（MIB）}：网络中被管理的每个设备上用于收集并储存管理信息的数据库，供网络管理系统获取。\\
\textbf{网络安全}：网络系统的硬件、软件及数据受到保护，不遭受偶然或恶意的破坏、更改、泄露，系统连续可靠正常地运行，网络服务不中断。\\
\textbf{机密性}：保证信息与信息系统不被非授权者获取与使用，主要防范措施是密码技术。\\
\textbf{完整性}：保证信息不被篡改或破坏，可通过密码技术提供完整性保护。\\
\textbf{可用性}：保证信息与信息系统可被授权者正常使用，主要防范措施是确保系统处于可靠的运行状态。
\end{keyterms}

\begin{contrast}{OSI 参考模型 / TCP/IP 协议族（每行依次为：分层·性质与地位·主要特点）}
\citem{OSI}{七层：物理层、数据链路层、网络层、运输层、会话层、表示层、应用层}{由 ISO 制定的概念性框架，为制定标准提供参考}{只有协议可以真正实现，设备需与 OSI 及有关协议一致才能互连}
\citem{TCP/IP}{四层：网络接口层、网络层、传输层、应用层}{Internet 实际使用的协议族，事实上的工业标准}{实现成熟、应用广泛，但分层不如 OSI 严格}
\end{contrast}

\begin{contrast}{语义 / 语法 / 同步（网络协议三要素；每行依次为：含义·解决的问题·举例）}
\citem{语义}{涉及用于协调与差错处理的控制信息}{规定通信双方“做什么”}{如差错控制、协调控制信息}
\citem{语法}{涉及数据及控制信息的格式、编码及信号电平}{规定通信双方“怎么做”}{如帧的格式、比特编码方式}
\citem{同步}{涉及速度匹配和排序等}{规定通信双方“何时做”}{如收发速度匹配、报文先后次序}
\end{contrast}

\begin{contrast}{TCP / UDP（每行依次为：连接方式·可靠性·适用场合）}
\citem{TCP}{传输控制协议，面向连接}{为应用程序提供可靠的通信连接}{适合一次传输大批数据、要求得到响应的应用，如文件传送、网页请求}
\citem{UDP}{用户数据报协议，无连接}{不对传送包进行可靠保证}{适合一次传输小量数据，可靠性由应用层负责，如实时音频视频}
\end{contrast}

\begin{contrast}{三网融合 / 物理合一（每行依次为：含义·融合层面·实质）}
\citem{三网融合}{技术功能趋于一致、业务范围趋于相同、互联互通、资源共享}{涉及技术、业务、行业、终端、网络五个层面的融合}{主要是高层业务应用的融合，不要求三大网络合并}
\citem{物理合一}{三大网络合并为一张物理网络}{是网络基础设施层面的合并}{并非三网融合的要求，三网融合不意味着物理合一}
\end{contrast}

\qtypename{一、单项选择题}{10 题，每题 3 分，共 30 分}
\q 在计算机网络的发展历程中，以美国国防部高级研究计划局开发的 ARPANET 为典型代表的是（\quad）
\choicesline{A. 诞生阶段}{B. 形成阶段}{C. 互联互通阶段}{D. 高速网络技术阶段}

\q 三网融合主要是指（\quad）
\choices{A. 三大网络的物理合一}{B. 高层业务应用的融合}
{C. 三家企业合并经营}{D. 三种终端设备的统一}

\q 网络协议的三要素是（\quad）
\choicesline{A. 语义、语法、同步}{B. 硬件、软件、数据}
{C. 输入、处理、输出}{D. 发送、接收、确认}

\q 计算机网络各层次及其协议的集合称为（\quad）
\choicesline{A. 网络拓扑}{B. 网络协议}{C. 网络体系结构}{D. 网络操作系统}

\q 在 OSI 七层模型中，第一个端到端（主机—主机）的层次是（\quad）
\choicesline{A. 网络层}{B. 运输层}{C. 会话层}{D. 表示层}

\q 在 OSI 七层模型中，负责数据语法表示变换以及数据压缩和加密的是（\quad）
\choicesline{A. 会话层}{B. 表示层}{C. 应用层}{D. 数据链路层}

\q 在 TCP/IP 协议族中，负责在主机和网络之间寻址和路由数据包的协议是（\quad）
\choicesline{A. IP}{B. TCP}{C. UDP}{D. HTTP}

\q 在传输层中，提供无连接通信、且不对传送包进行可靠保证的协议是（\quad）
\choicesline{A. TCP}{B. UDP}{C. IP}{D. FTP}

\q SNMP 在 TCP/IP 协议族中属于（\quad）协议。
\choicesline{A. 应用层}{B. 传输层}{C. 网络层}{D. 网络接口层}

\q 网络安全具有三个基本属性，即（\quad）
\choices{A. 机密性、完整性、可用性}{B. 开放性、共享性、透明性}
{C. 快速性、准确性、经济性}{D. 集中性、分布性、可靠性}

\qtypename{二、多项选择题}{3 题，每题 4 分，共 12 分；每题有两个或两个以上正确项}
\q 三网融合涉及下列哪些层面的融合（\quad）
\choicesline{A. 技术融合}{B. 业务融合}{C. 行业融合}{D. 终端融合}

\q 下列属于 OSI 七层参考模型层次的有（\quad）
\choicesline{A. 物理层}{B. 网络层}{C. 表示层}{D. 网络接口层}

\q 网络安全技术机制主要包括（\quad）
\choices{A. 加密和隐藏}{B. 认证和授权}
{C. 审计和定位}{D. 随意开放共享}

\qtypename{三、判断题}{6 题，每题 2 分，共 12 分；正确的画“√”，错误的画“×”并写出改正后的表述}
\q 三网融合意味着电信网、广播电视网和互联网合并成一张物理网络。\quad（\quad）\quad 改正：\fixblank

\q 现代移动通信已经历 1G、2G、3G、4G 四个发展阶段，其中 1G 是模拟、仅限语音的蜂窝电话标准。\quad（\quad）\quad 改正：\fixblank

\q OSI 参考模型是一个具体的网络实现方案，任何网络都必须按 OSI 实现。\quad（\quad）\quad 改正：\fixblank

\q 运输层是 OSI 模型中第一个端到端（主机—主机）的层次。\quad（\quad）\quad 改正：\fixblank

\q TCP 是无连接的传输协议，UDP 是面向连接的可靠协议。\quad（\quad）\quad 改正：\fixblank

\q 网络安全的三个基本属性是机密性、完整性和可用性。\quad（\quad）\quad 改正：\fixblank

\qtypename{四、配对题}{6 对，每对 1 分，共 6 分；把 A 栏与 B 栏对应内容配对}
\q 将下列网络层次、协议与其主要功能配对，把 B 栏字母填入括号。
\par\vspace{0.2em}
\noindent\begin{tabular}{@{}p{0.38\textwidth}@{\hspace{1.2em}}p{0.50\textwidth}@{}}
（1）物理层 \hfill（\quad） & A. 使原始比特流在物理媒体上传输 \\
（2）数据链路层 \hfill（\quad） & B. 把不可靠物理链路改造成对网络层无差错的数据链路 \\
（3）网络层 \hfill（\quad） & C. 分组传输、路由选择与网际互连 \\
（4）运输层 \hfill（\quad） & D. 提供端到端的透明数据传输 \\
（5）应用层 \hfill（\quad） & E. 为特定类型的网络应用提供访问 OSI 环境的手段 \\
（6）SNMP \hfill（\quad） & F. 基于 TCP/IP 协议族的网络管理标准 \\
\end{tabular}\par\vspace{0.3em}

\newpage
\qtypename{五、简答题}{4 题，每题 6 分，共 24 分}
\q 什么是三网融合？三网融合涉及哪几个层面的融合？其主要支撑技术有哪些？\anslines{2.4cm}

\q 简述 OSI 参考模型的七层结构及各层的主要功能。\anslines{2.4cm}

\q 什么是网络管理？简述网络管理的主要功能与 SNMP 的组成。\anslines{2.4cm}

\q 什么是网络安全？简述网络安全的三个基本属性与主要技术机制。\anslines{2.4cm}

\newpage
\qtypename{六、综合题（论述）}{1 题，共 16 分}
\q 某农业园区建设“智慧农业园区网络”：园区内布设传感器网络和摄像头，通过光纤接入互联网，同时接入有线电视和 IP 电话，向农户提供远程监控、视频培训与农产品电商服务，并与上级农业部门的信息系统互联，园区还要求保障数据安全。
请回答：
（1）结合 OSI 参考模型或 TCP/IP 协议族，说明数据从园区传感器传输到用户终端的大致过程；
（2）说明三网融合在该园区有哪些体现；
（3）说明应如何对该园区网络进行管理；
（4）从网络安全的角度，提出保障园区数据安全的技术与管理措施。\anslines{5cm}

\end{document}
